\documentclass[11pt,a4paper]{article}
\usepackage[utf8]{inputenc}
\usepackage[T1]{fontenc}
\usepackage{amsmath,amssymb,amsfonts,amsthm,mathtools}
\usepackage{graphicx}
\usepackage[margin=1in]{geometry}
\usepackage{fancyhdr}
\usepackage{enumitem}
\usepackage{multicol}
\usepackage{xcolor}
\usepackage{tikz}
\usepackage{subcaption}
\usepackage{kotex}
\usepackage{cancel}
\usepackage[colorlinks=true]{hyperref}
\usepackage{xcolor}
\newenvironment{solution}{%
    \par\noindent\color{blue}\textbf{Answer)}\par
}{%
    \par
}

\pagestyle{fancy}
\fancyhf{}
\renewcommand{\headrulewidth}{0.4pt}
\renewcommand{\footrulewidth}{0.4pt}

\lhead{\textbf{EN5232}}
\chead{\textbf{Climate Dynamics}}
\rhead{\textbf{2026-01}}

\cfoot{-~\thepage~-}

\newcounter{question}

\newcommand{\hwquestion}[2][]{ 
  \stepcounter{question} 
  \vspace{0.6cm} 
  \noindent 
  \ifx\relax#1\relax
  \else
    {\footnotesize(#1)}\\
  \fi
  \noindent \textbf{Q\thequestion. }#2 
  \vspace{0.2cm} 
}

\newtheorem{theorem}{Theorem}[section]

\definecolor{neongreen}{RGB}{57,255,20}
\definecolor{vividmagenta}{RGB}{255,0,255}
\definecolor{brightorange}{RGB}{255,128,0}
\definecolor{vividskublue}{RGB}{0,191,255}
\definecolor{brightgold}{RGB}{255,215,0}
\definecolor{brightlime}{RGB}{191,255,0}
\definecolor{vividpurple}{RGB}{159,0,255}
\definecolor{forestgreen}{RGB}{34,139,34}
\hypersetup{
    colorlinks=true,
    linkcolor=vividskublue,
    citecolor=brightlime,
}

\begin{document}
\providecommand{\hwtitle}{Midterm}
\providecommand{\hwduedate}{2026.04.24}
\providecommand{\studentid}{20251214}
\providecommand{\studentname}{김희수}

\begin{center}
    \Large~\textbf{\hwtitle}
\end{center}

\vspace{0.4cm}

\noindent
\begin{minipage}[t]{0.5\textwidth}
Due: \hwduedate
\end{minipage}
\begin{minipage}[t]{0.5\textwidth}
    \begin{flushright}
    Student ID: \studentid\\
    Name: \studentname
    \end{flushright}
\end{minipage}

\vspace{0.4cm}
\noindent\rule{\textwidth}{0.4pt}


\begin{center}
    \textbf{HW \#2} \\
    Due: April 15 (Wed)
\end{center}

\vspace{0.3cm}
\noindent\textbf{Answer the following questions}

\hwquestion[]{Why is the zonally averaged cross-equatorial overturning circulation, which can be depicted by zonally averaged meridional wind, stronger in JJA than in DJF?}
\begin{solution}

\end{solution}

\hwquestion[]{Why is it necessary to insert extra contours in the tropics as shown in the following Figure?}
\begin{figure}[h!]
    \centering
    \includegraphics[width=0.55\textwidth]{hw/images/hw2_q2_u200_Z200.png}
    \caption*{$\overline{u}_{200}$, $\overline{Z}_{200}$ for DJF and JJA.}
\end{figure}
\begin{solution}

\end{solution}

\newpage
\hwquestion[]{Prove that on average over the Hadley cell, the horizontal flow is down the pressure gradient.}
\begin{solution}

\end{solution}

\hwquestion[]{How would the overturning circulation depicted in Figure below be different in a laboratory experiment in which the heat sink is a cold top and right side boundary, rather than in the interior of the fluid?}
\begin{figure}[h!]
    \centering
    \includegraphics[width=0.45\textwidth]{hw/images/hw2_q4_lab_experiment.png}
    \caption*{Lab experiment: heated bottom, heated side.}
\end{figure}
\begin{solution}

\end{solution}

\newpage
\noindent\textbf{Energy Balance Model}

\hwquestion[]{A planet has the energy balance model shown below. In this case, there is no atmosphere. Red indicates incoming solar energy and black outgoing terrestrial radiation.}
\begin{figure}[h!]
    \centering
    \includegraphics[width=0.55\textwidth]{hw/images/hw2_q5_ebm_noatm.png}
    \caption*{Energy balance model without atmosphere.}
\end{figure}
\begin{enumerate}[label=(\alph*)]
    \item Set up the energy balance equation between incoming solar energy and outgoing radiation.
    \begin{solution}

    \end{solution}
    \item Compute the surface temperature of this planet, $T_S$, with the Stefan--Boltzmann constant $\sigma = 5.68 \times 10^{-8}\,\mathrm{W\,m^{-2}\,K^{-4}}$, the planetary albedo $\alpha_p$ as $0.3$, and solar constant $S_0 = 1360\,\mathrm{W\,m^{-2}}$.
    \begin{solution}

    \end{solution}
\end{enumerate}

\newpage
\hwquestion[]{Planet has the energy balance model shown below, which now has atmosphere, similar to the Earth.}
\begin{figure}[h!]
    \centering
    \includegraphics[width=0.55\textwidth]{hw/images/hw2_q6_ebm_atm.png}
    \caption*{Energy balance model with a single-layer atmosphere.}
\end{figure}
\begin{enumerate}[label=(\alph*)]
    \item In this case, \underline{atmosphere is completely transparent to solar radiation} while \underline{completely opaque to terrestrial radiation}. Set up the energy balance equation for atmosphere and surface.
    \begin{solution}

    \end{solution}
    \item Calculate the surface temperature ($T_S$), the atmosphere temperature ($T_A$), and the emission temperature ($T_E$).
    \begin{solution}

    \end{solution}
    \item Discuss the similarity and difference among these temperatures.
    \begin{solution}

    \end{solution}
    \item Now, let's introduce emissivity, $\epsilon$, which determines fraction of radiation that is emitted. The radiative flux emitted to space, $F$, through the top of the atmosphere is the sum of the radiative fluxes emitted by the top most layer $F_A$, and the transmitted radiative flux, $(1-\epsilon)F_S$, where $F_S$ is the upwelling radiative flux from the surface. Prove that in the limit as $\epsilon \to 0$, the temperature of the surface of a planet approaches the planet's equivalent blackbody temperature $T_E$ and the temperature of its atmosphere approaches $T_E / 2^{1/4}$.
    \begin{solution}

    \end{solution}
\end{enumerate}

\newpage
\noindent\textbf{Answer the following question after you read papers by E.\,Lorenz 1983 and I.\,Held 2019 (only up to Section 2).}

\hwquestion[]{What are missing and not logical in E.\,Halley's explanation on the trade wind?}
\begin{solution}

\end{solution}

\hwquestion[]{According to G.\,Hadley (1735), he explained the trade winds in terms of the angular momentum conservation. Please summarize it using your own words in 2--3 sentences.}
\begin{solution}

\end{solution}

\hwquestion[]{Concept of the angular momentum conservation is indeed consistent with the Coriolis force. If we average over longitude, the zonal momentum equation can be simplified in the following form:
$$\frac{\partial u}{\partial t} \cong f v$$
Using the idea of differential solar heating and the above equation instead of the angular momentum conservation as in 2), please write down the mechanism of the trade winds.}
\begin{solution}

\end{solution}

\hwquestion[]{What is the $\beta$-effect?}
\begin{solution}

\end{solution}

\newpage
\noindent\textbf{Please show the following}

\hwquestion[]{$v = \overline{v_Z} + \overline{v_E} + v'_Z + v'_E$ and $X = \overline{X_Z} + \overline{X_E} + X'_Z + X'_E$ as defined in the class. Show that
$$\overline{(vX)}_Z = \overline{v_Z}\,\overline{X_Z} + \overline{v'_Z X'_Z} + \overline{(\overline{v_E} \cdot \overline{X_E})}_Z + \overline{(v'_E X'_E)}_Z.$$}
\begin{solution}

\end{solution}

\hwquestion[]{Begin with the continuity equation in spherical and pressure coordinate
$$\frac{1}{a\cos\varphi}\left[\frac{\partial u}{\partial \lambda} + \frac{\partial}{\partial \varphi}(v\cos\varphi)\right] + \frac{\partial \omega}{\partial p} = 0$$}
\begin{enumerate}[label=(\alph*)]
    \item Show that $\int_0^{p_0} v_Z\,dp = 0$.
    \begin{solution}

    \end{solution}
    \item What is the physical meaning of this?
    \begin{solution}

    \end{solution}
\end{enumerate}

\end{document}
